\section{The next year starts from closed group equity}
\label{sec:worked-intragroup-carry}

Suppose the first year ends with the entire lot still inside the group. The
next report cannot eliminate that year's internal sale as if it had happened
again. The parent's legal margin has been closed into retained profit, and
the group's opening inventory already excludes it. The first task is therefore
to carry the two legal books and the consolidation adjustments into the new
year without changing the preceding report.

The legal close transfers each entity's income and expense balances into
retained profit. The group close also transfers the acquired-basis depreciation
and related tax effect into the subsidiary's retained result. NCI's first-year
profit is then moved out of owners' retained profit and added to the NCI
already recognised at acquisition. In the unsold-inventory example, group
opening equity remains $1521/4$: parent owners hold $2749/8$ and NCI holds
$293/8$. Carrying the original acquisition NCI alone would omit its first-year
profit; adding the whole closing NCI again would count part of it twice.

A separate opening worksheet records the account-by-account difference between
the closed legal books and those closed group accounts. It retains the
investment/equity elimination, acquired equipment basis, goodwill, inventory
margin and tax adjustments, with their original journal sources. This is a
carried balance, not another acquisition. Its complete accounts must reproduce
the first year's closing assets, liabilities, cash and both equity interests
before any second-year transaction is evaluated.

\subsection{Releasing a carried margin changes current cost of sales}

In the second year, let the subsidiary sell the full lot outside the group
for 40 cash, under the same supplied revenue-recognition premises. Its legal
cost is 30, but the group opening inventory is 20. The legal sale removes
30 from the subsidiary's inventory; the opening group adjustment still contains
a reduction of 10. Releasing that carried adjustment debits inventory 10 and
credits current cost of sales 10. The group therefore reports cost 20 and
zero closing inventory. Debiting opening retained profit a second time would
duplicate an adjustment already present in the opening worksheet.

Under the supplied tax assumptions, the subsidiary's current tax on its legal
profit is $5/2$, and the carried inventory deferred tax asset of $5/2$ reverses.
The inventory transaction consequently contributes after-tax group profit 15.
The subsidiary contributes $15/2$ after its current tax; the remaining $15/2$
is the parent-attributed release of the downstream margin, net of the deferred
tax reversal. These attribution premises remain explicit. An upstream seller
would require a different comparison.

\subsection{Other acquired assets continue to affect the year's result}

The equipment did not disappear at the inventory sale. A second-year comparison
that freezes it would omit a known accounting dependency. Continue the
synthetic example with separately supplied full-year legal depreciation 14,
additional acquired-basis depreciation 4 and a tax-liability release of 1.
The amounts extend the first year's half-year pattern, but are declared
schedules, not inferred useful lives or tax deductions. AASB 10 requires the
subsidiary's consolidated expenses to retain the acquisition basis; it does
not determine these particular amounts
\citep[para.~B88]{aasb10control2026}.

Equipment contributes a net expense of 17: legal depreciation 14 plus the
additional 4, less tax income 1. Combined with inventory profit 15, second-year
group profit is therefore $-2$. The subsidiary-attributed result is
$15/2-17=-19/2$, so ordinary NCI receives $-19/8$. Parent owners receive
three quarters of that result plus the parent-attributed margin release,
giving $3/8$. Their sum is $-2$. The negative NCI result is retained rather
than clipped \citep[para.~B94]{aasb10control2026}.

The closing group holds cash 310, equipment 63 and goodwill $135/4$, with no
remaining inventory or inventory deferred tax asset. Assets total $1627/4$.
Liabilities are $57/2$: the original external payable 20, acquired deferred
tax liability $7/2$, and unpaid current taxes of 5. Equity is consequently
$1513/4$, equal to opening equity $1521/4$ less the current loss of 2.
The earlier tax payable has not vanished, and no payment has been invented
to make the roll-forward balance.

The reference compares complete closed openings, the carried adjustment
worksheet, second-year legal and group journals, closing accounts and cash
origins. No-sale and partial-sale cases retain the appropriate inventory and
tax amounts while still recognising the supplied equipment expenses. Unknown
tax continuation retains the prior report and proposed request but selects
no second-year final report.\footnote{The full comparison and protected
inputs are retained in
\path{data/accounting-master/acquisition-intragroup-carry-20261027/}.
The pre-execution audit repaired the omitted equipment schedule and avoided
a duplicate debit to opening retained profit before running the candidate.}
This remains a one-lot downstream comparison with supplied continuing control,
ordinary rights and financial schedules. The next discriminating case must
place the seller inside the partially owned subsidiary and check which profit
adjustments then belong to NCI.
